\documentclass[10pt,a4paper]{amsart}
\usepackage[margin=19mm]{geometry}
\usepackage{amsmath,amssymb,mathtools}
\usepackage[scr=rsfs,cal=euler]{mathalfa}
\usepackage{xfrac}
\usepackage[unicode,hidelinks,bookmarks=false]{hyperref}
\newcommand{\ie}{i.e.\ }
\newcommand{\cf}{cf.\ }
\newcommand{\resp}{respectively\ }
\newcommand{\Subsection}{\subsection}
\providecommand{\nobreakdash}{\nobreak\hbox{-}}
\setlength{\emergencystretch}{2em}
\multlinegap0pt
\usepackage{bm}
\usepackage{bbm}
\usepackage{dsfont}
\usepackage{xspace}
\usepackage{cancel}
\usepackage{mathtools}
\usepackage{cleveref}
\usepackage{amsthm}
\makeatletter
\@ifundefined{proposition}{\newtheorem{proposition}{Proposition}}{}
\@ifundefined{lemma}{\newtheorem{lemma}[proposition]{Lemma}}{}
\makeatother
\DeclareMathOperator{\Frob}{Frob}
\DeclareMathOperator{\Per}{Per}
\DeclareMathOperator{\car}{char}
\DeclareMathOperator{\tr}{tr}
\def\FF{\mathbb F}
\def\PP{\mathbb P}
\def\abs#1{\left\vert{#1}\right\vert}
\def\set#1{{\def\st{\;:\;}\left\{#1\right\}}}
\title[Proportion of periodic points in reduction of polynomials]{Proportion of periodic points in reduction of polynomials}
\author{Santiago Radi}
\date{Source: arXiv:2603.21620v2 (CC BY 4.0)}
\begin{document}
\maketitle

\noindent\emph{Source and licence.} Proportion of periodic points in reduction of polynomials. Version arXiv:2603.21620v2 (CC BY 4.0), distributed under the Creative Commons Attribution 4.0 International licence, as recorded in the arXiv licence metadata for this exact version and shown on the abstract page https://arxiv.org/abs/2603.21620v2. Full rights notice: licenses/arxiv-2603.21620-CC-BY-4.0.txt.\par\smallskip

\noindent\textit{Editorial scope.} Counted unit: the Proposition whose source label is \texttt{proposition: Perinf f with a not in K} in the article (source0.tex lines 715--721, source label \texttt{proposition: Perinf f with a not in K}), delivered together with the article's own complete original proof of it (source0.tex lines 723--808, one \texttt{proof} environment of 86 lines). No mathematics is written, repaired or re-derived: statement and proof are the article's own text line for line. Required context retained uncounted: lemma: Per -Td and Per Td (lines 516--520); lemma: Per Td (lines 587--594); lemma: number of elements with order coprime to a number (lines 617--624); lemma: number of Per for Td (lines 636--640). Every definition, display and label of those lines is retained unchanged. Omitted from this package and not redistributed: 1--515, 521--586, 595--616, 625--635, 641--714, 722--722, 809--1015. Nothing inside a retained excerpt is omitted or altered: every retained body is delivered exactly as the article's lines are written, so no omission receipt of any kind accompanies this package. Citations: the retained excerpts carry no citation command at all, so no bibliography entry of the article is transcribed and no reference is redistributed. Reference adaptation: none was needed and none was made; every reference inside the delivered unit resolves inside it, so no unresolved reference or citation is produced. Printed slips kept literal: no printed word, symbol, hypothesis or number of the counted statement or of its proof was altered. Layout and class: the article's own class is \texttt{amsart} with packages including inputenc, fontenc, textcomp, amsfonts, amsthm, amssymb, pst-node, xy, mathtools, chngcntr, thmtools, mathrsfs, tcolorbox, listings; the wrapper is the lane's pinned \texttt{amsart} editorial wrapper at 10pt on A4, which restates the theorem-like environments the retained text needs and adds the article's own macro definitions that the retained text needs (\texttt{\textbackslash FF}, \texttt{\textbackslash Frob}, \texttt{\textbackslash PP}, \texttt{\textbackslash Per}, \texttt{\textbackslash abs}, \texttt{\textbackslash car}, \texttt{\textbackslash set}, \texttt{\textbackslash tr}), with extra wrapper packages (bm, bbm, dsfont, xspace, cancel, mathtools, cleveref). The delivered numbering is the unit's own. Assets: no retained span contains a figure environment, a graphics-inclusion command, a PDF-page inclusion, a table or a TikZ picture; no image file is copied into this package, the package declares no asset dependency and no image file is redistributed with it. Licence: version arXiv:2603.21620v2 of this article is distributed under the Creative Commons Attribution 4.0 International licence, as recorded in the arXiv licence metadata for this exact version and shown on the abstract page https://arxiv.org/abs/2603.21620v2. A full rights notice is delivered at licenses/arxiv-2603.21620-CC-BY-4.0.txt. Characters: every retained excerpt is pure ASCII, so the delivered engine, XeLaTeX with its default Latin Modern text font, reports no missing character.
\begin{lemma}
\label{lemma: Per -Td and Per Td}
Given $d \geq 2$ and a field $F$, then 
$$\Per(-T_d, F) =  (-1)^{d+1}\Per(T_d, F).$$ 
\end{lemma}

\begin{lemma}
Let $d \geq 2$, $F$ a finite field with $\car(F) \neq 2$, $F_2$ the extension of degree $2$ of $F$ and let $\abs{ \cdot}$ denote the order of an element in the group $F_2^\times$. Then 
\begin{align*}
\Per(T_d, F) = \bigcup_{\gamma \in \set{-1, 1}} \set{\pi(\beta) \in F: \gcd(\abs{\beta},d) = 1 \text{ and } \abs{\beta} \mid (\abs{F} + \gamma)} 
\end{align*}
and the sets involved in the union intersects at $1$ if $d$ is even or at $\pm 1$ if $d$ is odd.
\label{lemma: Per Td}
\end{lemma}

\begin{lemma}
\label{lemma: number of elements with order coprime to a number}
Let $G$ be a cyclic finite group of order $n$. Let $m \geq 1$ such that $m \mid n$ and let $b \geq 1$. Define $$H := \set{g \in G: g^m = 1}.$$
Then, 
\begin{align*}
\# \set{h \in H: \gcd(\abs{h},b) = 1} = r(m,b).
\end{align*}
\end{lemma}

\begin{lemma}
\label{lemma: number of Per for Td}
Let $d \geq 2$ and $F$ a finite field with $\car(F) \neq 2$. Then 
$$\# \Per(\pm T_d, F) = \frac{r(\abs{F}-1,d) + r(\abs{F}+1,d) }{2}.$$
\end{lemma}

\begin{proposition}
\label{proposition: Perinf f with a not in K}
Let $K$ be a number field, an odd number $d \geq 2$ and consider $f = L^{-1} \circ (\pm T_d) \circ L \in K[x]$ with $L(x) = ax$ and $a \notin K$. Then,  
\begin{align*}
\Per_{\inf}(f,K) = \frac{1}{2} \liminf_{N(\mathfrak{p}) \rightarrow +\infty} \frac{r(N(\mathfrak{p})-1,d) + r(N(\mathfrak{p})+1,d)}{N(\mathfrak{p})}.
\end{align*}
\end{proposition}

\begin{proof}
Call $K' = K(a)$. By hypothesis $K'/K$ is an extension of degree $2$. Let $\mathfrak{p}$ be a prime ideal such that $\car(\FF_\mathfrak{p}) \neq 2$ and $\mathfrak{p}$ does not ramify at $K'$. Notice that we are excluding only finitely many prime ideals. Let $\mathfrak{p}'$ be a prime ideal in $\mathcal{O}_{K'}$ such that $\mathfrak{p}' \mid \mathfrak{p}$.

As $f$ and $\pm T_d$ are linearly conjugate over $K'$, then $L$ induces a bijection
\begin{align}
\label{equation: bijection Per f and Per Td d odd a not in K}
\Per(f_{\mathfrak{p}'}, \FF_{\mathfrak{p}'}) &\rightarrow  \Per((\pm T_d)_{\mathfrak{p}'}, \FF_{\mathfrak{p}'}) \\
\notag \alpha &\mapsto a \alpha.
\end{align}
Moreover, since $d$ is odd by \cref{lemma: Per -Td and Per Td} we have $\Per((\pm T_d)_{\mathfrak{p}'}, \FF_{\mathfrak{p}'}) = \Per((T_d)_{\mathfrak{p}'}, \FF_{\mathfrak{p}'})$.

We have two cases for $\mathfrak{p}$. If $\mathfrak{p}$ is split, then $\FF_{\mathfrak{p}'} = \FF_\mathfrak{p}$ and therefore by \cref{lemma: number of Per for Td},
\begin{align}
\# \Per(f_\mathfrak{p}, \FF_\mathfrak{p}) = \frac{r(N(\mathfrak{p})-1,d) + r(N(\mathfrak{p})+1,d)}{2}.
\label{equation: Per f in split case}
\end{align}

If $\mathfrak{p}$ is inert, then $\FF_{\mathfrak{p}'}/\FF_\mathfrak{p}$ is an extension of degree $2$ and consequently the Frobenius map sends $a$ to $-a$ in $\FF_\mathfrak{p'}$.

Recall the trace is (in this case) the surjective additive group homomorphism
\begin{align*}
\tr: \FF_{\mathfrak{p}'} &\rightarrow \FF_\mathfrak{p}, \\
y &\mapsto y + \Frob_\mathfrak{p}(y).
\end{align*}
As it is surjective we have $$\# \ker(\tr) = N(\mathfrak{p}).$$

We claim that
\begin{align*}
a \FF_\mathfrak{p} = \ker(\tr).
\end{align*}

Indeed, if $y\in a \FF_\mathfrak{p}$ then $y = ax$ for some $x \in \FF_\mathfrak{p}$. Therefore 
\begin{align*}
\tr(y) = ax -ax^{N(\mathfrak{p})} = 0
\end{align*}
as $x \in \FF_\mathfrak{p}$. This proves that $a\FF_\mathfrak{p} \subseteq \ker(\tr)$, but since they have the same cardinality, then $a\FF_\mathfrak{p} = \ker(\tr)$.

Let $F_2$ be the extension of degree $2$ of $\FF_{\mathfrak{p}'}$ and recall the map $\pi: \PP^1(F_2) \rightarrow \PP^1(F_2)$. Then, if $\beta \in F_2^\times$,
\begin{align*}
\tr(\pi(\beta)) = \beta + \beta^{-1} + \beta^{N(\mathfrak{p})} + \beta^{-N(\mathfrak{p})} = \frac{(\beta^{N(\mathfrak{p})-1} + 1)(\beta^{N(\mathfrak{p})+1} + 1)}{\beta^{N(\mathfrak{p})}}.
\end{align*}

We are interested in calculating $\# \Per(f_\mathfrak{p}, \FF_\mathfrak{p})$. Using the bijection in \cref{equation: bijection Per f and Per Td d odd a not in K}, we obtain that 
\begin{align*}
\# \Per(f_\mathfrak{p}, \FF_\mathfrak{p}) = \#(a \FF_\mathfrak{p} \cap \Per((T_d)_{\mathfrak{p}'}, \FF_{\mathfrak{p}'})).
\end{align*}

Using \cref{lemma: Per Td}, we reduce the previous problem to calculate the number of elements of
\begin{align*}
V_\gamma := \set{\pi(\beta) \in \FF_{\mathfrak{p}'}: \tr(\pi(\beta)) = 0, \,\, \gcd(\abs{\beta},d) = 1 \text{ and } \abs{\beta} \mid (N(\mathfrak{p})^2+\gamma)} 
\end{align*}
for $\gamma \in \set{-1,1}$.

Start with $\gamma = 1$. We claim that 
$$\set{\beta \in F_2^\times: \tr(\pi(\beta)) = 0 \text{ and } \abs{\beta} \mid (N(\mathfrak{p})^2+1)} = \emptyset,$$ which in particular implies that $V_1 = \emptyset$. Indeed, if $\tr(\pi(\beta)) = 0$ then $\beta^{N(\mathfrak{p}) \pm 1} = -1$ for either $+$ or $-$ in the exponent. Therefore $$\abs{\beta} \mid 2(N(\mathfrak{p})\pm 1).$$
Notice that as $N(\mathfrak{p})$ is odd, the previous formula implies that $$\abs{\beta} \mid (N(\mathfrak{p})^2 - 1).$$
On the other hand if $\abs{\beta} \mid (N(\mathfrak{p})^2 + 1)$, subtracting we conclude that $\abs{\beta} \mid 2$ and therefore $\beta = \pm 1$. However $\tr(\pi(\pm 1)) \neq 0$.

If $\gamma = -1$ consider 
$$H_{-1} := \set{\beta \in F_2^\times: \tr(\pi(\beta)) = 0 \text{ and } \abs{\beta} \mid (N(\mathfrak{p})^2-1)}.$$
By the calculation done before if $\tr(\pi(\beta)) = 0$ then $\abs{\beta} \mid (N(\mathfrak{p})^2-1)$, so
$$H_{-1} = \set{\beta \in F_2^\times: \tr(\pi(\beta)) = 0}.$$
Let $$G_c := \set{\beta \in F_2^\times: \beta^c = 1}.$$ 

Then
$$H_{-1} = \bigcup_{\gamma \in \set{-1,1}} (G_{2(N(\mathfrak{p})+\gamma)} \setminus G_{N(\mathfrak{p})+\gamma})$$ as $\tr(\pi(\beta)) = 0$ if and only if $\beta^{N(\mathfrak{p}) + \gamma} = -1$ for either $\gamma = 1$ or $\gamma = -1$.

Moreover $$\beta \in \bigcap_{\gamma \in \set{-1,1}} (G_{2(N(\mathfrak{p})+\gamma)} \setminus G_{N(\mathfrak{p})+\gamma})$$ if and only if $\abs{\beta} = 4$. 

Notice that as $\pm 1 \notin H_{-1}$, if $\pi(\beta) \in V_{-1}$ then $\pi(\beta)$ has two different preimages in $H_{-1}$. Using this observation and \cref{lemma: number of elements with order coprime to a number} we obtain
\begin{align*}
\# V_{-1} = \frac{1}{2} \left(\sum_{\gamma \in \set{-1,1}} (r(2(N(\mathfrak{p}) + \gamma),d) - r(N(\mathfrak{p}) + \gamma,d)) \right) - \frac{\#\set{\beta \in F_2^\times: \abs{\beta} = 4}}{2}.
\end{align*}
As $d$ is odd then $r(2(N(\mathfrak{p}) + \gamma),d) = 2 r(N(\mathfrak{p}) + \gamma,d)$ and therefore
\begin{align*}
 \# V_{-1} = \frac{1}{2} \left( \sum_{\gamma \in \set{-1,1}} r(N(\mathfrak{p}) + \gamma,d) \right) - \frac{\#\set{\beta \in F_2^\times: \abs{\beta} = 4}}{2}.
\end{align*}
As $V_1$ was empty, then $\# \Per(f_\mathfrak{p}, \FF_\mathfrak{p}) = \# V_{-1}$ allowing us to conclude
\begin{align}
\# \Per(f_\mathfrak{p}, \FF_\mathfrak{p}) = \frac{1}{2} \left( \sum_{\gamma \in \set{-1,1}} r(N(\mathfrak{p}) + \gamma,d) \right) - \frac{\#\set{\beta \in F_2^\times: \abs{\beta} = 4}}{2}.
\label{equation: Per f in inert case}
\end{align}

Dividing \cref{equation: Per f in split case} and \cref{equation: Per f in inert case} by $N(\mathfrak{p})$ and taking limit, we obtain
$$\Per_{\inf}(f,K) = \frac{1}{2} \liminf_{N(\mathfrak{p}) \rightarrow +\infty} \frac{r(N(\mathfrak{p})-1,d) + r(N(\mathfrak{p})+1,d)}{N(\mathfrak{p})}$$ as desired.
\end{proof}

\end{document}
